\section{Mục tiêu, phạm vi và môi trường triển khai}
\subsection{Yêu cầu bài thực hành}
Bài thực hành tích hợp định vị trên bản đồ có sẵn, điều hướng tự động và nhận diện ảnh RGB trên TurtleBot 4 mô phỏng. Bản đồ đã xây dựng ở bài 04 được sử dụng lại; robot tự đi qua ba mục tiêu có thứ tự, trong khi cửa sổ camera hiển thị kết quả YOLO cập nhật theo thời gian. Hai luồng điều hướng và nhận diện chạy đồng thời qua ROS 2.
\begin{table}[H]\centering\small
\caption{Yêu cầu và tiêu chí quan sát được}\label{tab:requirements}
\begin{tabularx}{\linewidth}{@{}P{3.5cm}Y@{}}\toprule
Nội dung&Tiêu chí hoàn thành\\\midrule
Nạp map và localization&Map cũ xuất hiện trong RViz2; AMCL hoạt động; TF nối từ map tới robot; LiDAR phù hợp với tường trên map.\\
Điều hướng tuần tự&Code gửi Goal\_1, Goal\_2, Goal\_3; mỗi action có kết quả SUCCEEDED trước khi gửi mục tiêu kế tiếp.\\
Nhận diện song song&Camera thật của simulator được xử lý thành bbox, lớp, confidence; có ảnh annotated trong lúc odometry cho thấy robot chuyển động.\\\bottomrule
\end{tabularx}\end{table}
\subsection{Môi trường và khả năng tái lập}
\begin{table}[H]\centering\small
\caption{Các thành phần phần mềm của workspace}
\begin{tabularx}{\linewidth}{@{}P{4cm}Y@{}}\toprule
Thành phần&Phiên bản/cấu hình\\\midrule
Hệ điều hành; ROS&Ubuntu 22.04; ROS 2 Humble; Python 3.10.12\\
Simulator; robot&Ignition Gazebo Fortress 6.18.0; world maze; TurtleBot 4 standard\\
Điều hướng&\code{turtlebot4_navigation}, Nav2, AMCL; tham số chính thức của bản Humble đang cài\\
Nhận diện hiện tại&Ultralytics 8.3.221; YOLO11s; ảnh đích 640; lọc sáu lớp COCO\\
Thư viện ảnh/tensor&PyTorch 2.5.1+cu121; torchvision 0.20.1+cu121; OpenCV 4.10.0.84; NumPy 1.26.4\\
GPU trong log kiểm chứng&NVIDIA GeForce RTX 3060 Laptop, 6144 MiB; CUDA runtime 12.1\\
Mạng ROS của bài&\code{ROS_DOMAIN_ID=64}; \code{ROS_LOCALHOST_ONLY=1}\\\bottomrule
\end{tabularx}\end{table}
\note{Báo cáo đối chiếu mã nguồn nền \code{a90f949} và kiểm tra package ngày 20/10/2026. Số liệu điều hướng có timestamp ngày 15/10/2026 được ghi rõ là lịch sử. Cấu hình hiện tại có 8 vật thể, YOLO11s và bộ lọc mới; không sử dụng số liệu của phiên bản YOLO cũ để khẳng định độ chính xác của phiên bản mới.}
\noindent Tài nguyên nộp bài: \ghlink{}; \vidlink{}. Các liên kết dùng nhãn ngắn và có thể bấm trực tiếp trong PDF.
\clearpage

\section{Cấu trúc chương trình và kiến trúc tổng thể}
\subsection{Phân chia trách nhiệm trong mã nguồn}
\begin{table}[H]\centering\small
\caption{Các tệp chính của package và vai trò}\label{tab:files}
\begin{tabularx}{\linewidth}{@{}P{5.25cm}Y@{}}\toprule
Tệp/thư mục&Trách nhiệm chính\\\midrule
\code{launch/bringup.launch.py}&Gọi launch localization/Nav2 của bài 04, RViz bài 05; có tùy chọn bật YOLO và mission. Mặc định chưa tự gửi goal.\\
\code{launch/yolo.launch.py}&Chạy riêng nhận diện; khai báo device, weights, ngưỡng và cửa sổ.\\
\code{sequential_navigation.py}&Action client bất đồng bộ, readiness, feedback, timeout, retry, cancel, dwell, log JSONL.\\
\code{configuration.py}&Kiểm tra waypoint, map hash, giá trị hữu hạn; đổi yaw sang quaternion; kiểm tra vùng trống.\\
\code{initialize_pose.py}&Khởi tạo AMCL trước mission từ pose đã xác thực hoặc nguồn mô phỏng được chọn tường minh.\\
\code{capture_waypoints.py}&Nhận điểm từ RViz, bổ sung yaw, kiểm tra map và lưu YAML.\\
\code{yolo_node.py}&Nhận ảnh mới nhất, suy luận trên worker, lọc bbox, xuất ảnh/JSON/status và cửa sổ OpenCV.\\
\code{detection_filter.py}&Xác nhận bbox cùng lớp qua các frame liền nhau bằng IoU; xóa trạng thái khi dữ liệu không hợp lệ/cũ.\\
\code{spawn_objects.py},\newline\code{scene_layout.py}&Chuẩn bị asset SDF; kiểm tra footprint với map và tường maze; thêm/thay model vào simulator.\\
\code{config/}, \code{maps/}&Waypoint, pose, YAML YOLO, vật thể, RViz; bản sao map từ bài 04.\\
\code{scripts/}, \code{validation/}&Môi trường Python, download, kiểm tra scene, observer; các snapshot kiểm chứng và ảnh thực tế.\\\bottomrule
\end{tabularx}\end{table}
Các module Python trong bảng nằm dưới \code{turtlebot4_lab05/}. Năm executable được khai báo trong \code{setup.py}: \code{sequential_navigation}, \code{yolo_node}, \code{initialize_pose}, \code{capture_waypoints} và \code{spawn_demo_objects}. Các script tiện ích được chạy trực tiếp bằng Python hoặc source trong shell.

Bài 05 tái sử dụng \code{turtlebot4_lab} để khởi động simulator/bridge và gọi localization/navigation chính thức. Cách này duy trì cùng world, robot, topic và map như bài trước, đồng thời tập trung phần mới vào mission và thị giác.
\clearpage

\subsection{Hai luồng xử lý chạy đồng thời}
\begin{figure}[H]\centering
\begin{tikzpicture}[node distance=9mm]
\node[block,text width=13.4cm] (sim) at (0,0) {\textbf{Gazebo Fortress -- maze -- TurtleBot 4 standard}\\LiDAR, odometry, RGB; thực thi chuyển động qua bridge/Create 3};
\node[block,text width=6cm] (amcl) at (-3.8,-2) {\textbf{Localization}\\Map server + AMCL + TF};
\node[block,text width=6cm,fill=softpurple] (rgb) at (3.8,-2) {\textbf{Camera RGB}\\OAK-D $\rightarrow$ cv\_bridge};
\node[block,text width=6cm,below=12mm of amcl] (nav) {\textbf{Nav2}\\BT navigator, NavFn, DWB, costmaps};
\node[block,text width=6cm,below=12mm of rgb,fill=softpurple] (yolo) {\textbf{YOLO worker}\\CUDA/CPU, ngưỡng lớp, xác nhận bbox};
\node[block,text width=6cm,below=12mm of nav,fill=softgreen] (seq) {\textbf{Sequential navigation}\\Waypoint $\rightarrow$ action; chờ kết quả};
\node[block,text width=6cm,below=12mm of yolo,fill=softpurple] (out) {\textbf{Quan sát nhận diện}\\Ảnh annotated, detections, cửa sổ};
\draw[arrow](sim.south -| amcl.north)--node[label]{scan, odom, TF}(amcl.north);
\draw[arrow](sim.south -| rgb.north)--node[label]{RGB image}(rgb.north);
\draw[arrow](amcl)--node[label]{pose trong map}(nav);
\draw[arrow](rgb)--(yolo);
\draw[arrow](yolo)--(out);
\draw[arrow](seq.north)--node[label]{NavigateToPose}(nav.south);
\draw[arrow](yolo.west)--++(-.45,0)|-node[label,pos=.75,above]{status}(seq.east);
\end{tikzpicture}
\caption[Kiến trúc tổng thể của hệ thống]{Kiến trúc chức năng. Readiness của mission dùng status YOLO; kết quả bbox không đi vào bộ điều khiển Nav2. Luồng vận tốc do Nav2 phát về robot được mô tả ở Hình~\ref{fig:navflow}.}\label{fig:architecture}
\end{figure}
Trong mã thực tế, status YOLO được gửi tới node mission. Mũi tên status trong sơ đồ biểu diễn điều kiện cho phép bắt đầu điều hướng; thông tin nhận diện không tạo vật cản trong costmap. Các quan sát vật cản dùng LiDAR qua \code{/scan}.

Một callback camera chỉ lưu ảnh mới nhất và đánh thức worker. Suy luận nặng không chạy trong callback nhận ảnh. Node mission dùng action bất đồng bộ nên vẫn xử lý được feedback, heartbeat và tín hiệu hủy trong khi Nav2 chạy.

\note{Chuyển động là vòng kín: cảm biến $\rightarrow$ AMCL/costmap $\rightarrow$ Nav2 $\rightarrow$ vận tốc $\rightarrow$ robot. Nhận diện là nhánh quan sát song song: camera $\rightarrow$ YOLO $\rightarrow$ hiển thị. Không có lệnh vận tốc được phát từ hai node Python của bài 05.}
\clearpage

\section{Package, node và giao diện ROS 2}
\subsection{Các node/package theo từng chức năng}
\begin{table}[H]\centering\small
\caption{Các thành phần tham gia pipeline}\label{tab:nodes}
\begin{tabularx}{\linewidth}{@{}P{3.1cm}P{5cm}Y@{}}\toprule
Nhóm&Node/package tiêu biểu&Vai trò\\\midrule
Mô phỏng&\code{turtlebot4_ignition_bringup}, bridge, Create 3&World/robot; chuyển dữ liệu Ignition--ROS; thực thi vận tốc và undock.\\
Map&\code{nav2_map_server}: \code{map_server}&Đọc YAML và PGM, publish OccupancyGrid đã lưu.\\
Định vị&\code{nav2_amcl}: \code{amcl}&Kết hợp map, scan và odometry; ước lượng pose và TF map--odom.\\
Lập kế hoạch&\code{planner_server}, \code{nav2_navfn_planner}&Tính đường toàn cục trên global costmap.\\
Điều khiển&\code{controller_server}, \code{dwb_core}&Chọn vận tốc bám đường và tránh vật cản cục bộ.\\
Điều phối Nav2&\code{bt_navigator}, \code{behavior_server}&Thực thi Behavior Tree; replan, follow path, recovery.\\
Costmap/vận tốc&\code{nav2_costmap_2d}, \code{velocity_smoother}&Cập nhật lớp static/obstacle/inflation; làm mượt vận tốc.\\
Lifecycle&\code{lifecycle_manager_*}&Configure/activate các node quản lý lifecycle; mission đọc GetState.\\
Mission bài 05&\code{lab05_sequential_navigation}&Gửi ba NavigateToPose; ghi kết quả từng UUID.\\
Thị giác bài 05&\code{lab05_yolo}, \code{cv_bridge}, Ultralytics&Chuyển RGB sang BGR; inference; xuất ảnh annotated và JSON.\\
Giao diện&\code{rviz2}, OpenCV HighGUI&Map, scan, TF, path, costmap và cửa sổ camera.\\\bottomrule
\end{tabularx}\end{table}
YOLO/PyTorch là thư viện Python được node sử dụng, không phải các node ROS độc lập. \code{nav2_simple_commander} có trong dependency, nhưng mission này trực tiếp tạo \code{rclpy.action.ActionClient}; không dùng \code{BasicNavigator} hay action FollowWaypoints để chạy chuỗi mục tiêu.
\subsection{Lifecycle và điều kiện active}
Mission gọi \code{/<node>/get_state} kiểu \code{lifecycle_msgs/srv/GetState}. Các node được kiểm tra gồm AMCL, map server, planner, controller, BT navigator, behavior server, velocity smoother và hai costmap. State id 3 là active. Tồn tại process hoặc xuất hiện trong \code{ros2 node list} chưa đủ để khẳng định pipeline đã sẵn sàng.
\clearpage

\subsection{Các topic chính và ý nghĩa dữ liệu}
\begin{table}[H]\centering\footnotesize
\caption{Topic, kiểu message và hướng truyền dữ liệu}\label{tab:topics}
\begin{tabularx}{\linewidth}{@{}P{4.3cm}P{4.2cm}Y@{}}\toprule
Topic&Kiểu message&Nguồn $\rightarrow$ nơi sử dụng\\\midrule
\code{/map}&\code{nav_msgs/msg/OccupancyGrid}&Map server $\rightarrow$ AMCL, static layer, RViz, mission.\\
\code{/scan}&\code{sensor_msgs/msg/LaserScan}&LiDAR/bridge $\rightarrow$ AMCL, costmap, RViz, readiness.\\
\code{/odom}&\code{nav_msgs/msg/Odometry}&Odometry Create 3 $\rightarrow$ Nav2, observer.\\
\code{/tf}, \code{/tf_static}&\code{tf2_msgs/msg/TFMessage}&AMCL, robot publishers $\rightarrow$ biến đổi tọa độ.\\
\code{/clock}&\code{rosgraph_msgs/msg/Clock}&Simulator $\rightarrow$ ROS time của các node.\\
\code{/initialpose}&\code{geometry_msgs/msg/}\code{PoseWithCovarianceStamped}&RViz/helper $\rightarrow$ khởi tạo phân bố AMCL.\\
\code{/amcl_pose}&\code{geometry_msgs/msg/}\code{PoseWithCovarianceStamped}&AMCL $\rightarrow$ RViz, mission, observer.\\
\code{/plan}, \code{/local_plan}&\code{nav_msgs/msg/Path}&Planner/DWB $\rightarrow$ RViz và quan sát đường đi.\\
\code{/cmd_vel}&\code{geometry_msgs/msg/Twist}&Nav2/smoother $\rightarrow$ bridge/Create 3.\\
\code{/oakd/rgb/preview/}\code{image_raw}&\code{sensor_msgs/msg/Image}&Camera RGB $\rightarrow$ YOLO, observer.\\
\code{/oakd/rgb/preview/}\code{camera_info}&\code{sensor_msgs/msg/CameraInfo}&Thông tin nội tại; pipeline bbox 2D không cần dùng trực tiếp.\\
\code{/yolo/annotated_image}&\code{sensor_msgs/msg/Image}&YOLO $\rightarrow$ RViz/observer; giữ header ảnh vào.\\
\code{/yolo/detections}&\code{std_msgs/msg/String} (JSON)&BBox, class, confidence, stamp và device.\\
\code{/yolo/status}&\code{std_msgs/msg/String} (JSON)&Readiness, freshness, số frame, thiết bị model.\\
\code{/lab05/navigation_status}&\code{std_msgs/msg/String} (JSON)&Trạng thái mission, UUID, kết quả, hoàn thành.\\
\code{/clicked_point}&\code{geometry_msgs/msg/PointStamped}&Publish Point $\rightarrow$ capture waypoint; không có yaw.\\
\code{/sim_ground_truth_pose}&\code{nav_msgs/msg/Odometry}&Đối chiếu mô phỏng; tùy chọn khởi tạo trước mission.\\
\code{/hazard_detection}&\code{irobot_create_msgs/msg/}\code{HazardDetectionVector}&Observer ghi các cảnh báo từ Create 3.\\\bottomrule
\end{tabularx}\end{table}
\noindent\textbf{QoS:} map, AMCL và status được subscribe bằng Transient Local để nhận trạng thái gần nhất khi node khởi động muộn. RGB và ảnh annotated dùng Best Effort, depth 1; mục tiêu là xử lý frame gần hiện tại, chấp nhận bỏ frame khi tải cao. QoS không phù hợp có thể khiến topic tồn tại nhưng subscriber không nhận dữ liệu.
\clearpage

\subsection{Action và service trong quá trình điều hướng}
\begin{table}[H]\centering\small
\caption{Các action/service chính}
\begin{tabularx}{\linewidth}{@{}P{4cm}P{4.3cm}Y@{}}\toprule
Tên&Kiểu&Vai trò\\\midrule
\code{/navigate_to_pose}&\code{nav2_msgs/action/}\code{NavigateToPose}&Mission gửi một PoseStamped; BT navigator nhận và điều phối.\\
\code{/compute_path_to_pose}&\code{nav2_msgs/action/}\code{ComputePathToPose}&BT gọi planner để tính global path.\\
\code{/follow_path}&\code{nav2_msgs/action/FollowPath}&BT yêu cầu controller bám path.\\
\code{/spin}, \code{/backup}, \code{/wait}&Action của \code{nav2_msgs}&Recovery tùy nhánh của Behavior Tree.\\
\code{/undock}&\code{irobot_create_msgs/action/Undock}&Đưa robot rời dock trước khi demo.\\
\code{/request_nomotion_update}&\code{std_srvs/srv/Empty}&Yêu cầu AMCL cập nhật khi robot đứng yên.\\
\code{/<node>/get_state}&\code{lifecycle_msgs/srv/GetState}&Mission kiểm tra state active.\\\bottomrule
\end{tabularx}\end{table}
\textbf{Vì sao dùng action?} Một lượt di chuyển có thể kéo dài hàng chục giây, cần phản hồi tiến độ và có thể hủy. Action cung cấp goal acceptance, UUID, feedback, result và cancel. Topic phù hợp luồng cảm biến liên tục; service phù hợp yêu cầu ngắn như đọc lifecycle hoặc no-motion update.

\code{NavigateToPose.Goal} gồm \code{pose} và \code{behavior_tree}. Code gán pose; chuỗi behavior tree để trống nên Nav2 dùng cây mặc định. Feedback của bản Humble có current pose, navigation time, estimated remaining time, recoveries và distance remaining. Code ghi khoảng cách còn lại và số recovery mỗi khoảng hai giây.

\textbf{Kết quả thành công đọc ở đâu?} Trường result của action Humble chỉ chứa Empty; status kết thúc nằm trong wrapper trả về bởi \code{get_result_async()}. Code đọc \code{response.status}: 4 = SUCCEEDED, 5 = CANCELED, 6 = ABORTED. Accepted chỉ cho biết server nhận yêu cầu, chưa chứng minh robot đã tới đích.

\textbf{Namespace của DWB:} RViz trong cấu hình này đọc \code{/local_plan}; khi đổi namespace/plugin cần kiểm tra lại bằng \code{ros2 topic list}. Topic path để hiển thị không thay thế trạng thái action làm tiêu chí hoàn thành mission.
